\documentclass[10pt,a4paper]{amsart}
\usepackage[margin=19mm]{geometry}
\usepackage{amsmath,amssymb,mathtools}
\usepackage[scr=rsfs,cal=euler]{mathalfa}
\usepackage{xfrac}
\usepackage[unicode,hidelinks,bookmarks=false]{hyperref}
\newcommand{\ie}{i.e.\ }
\newcommand{\cf}{cf.\ }
\newcommand{\resp}{respectively\ }
\newcommand{\Subsection}{\subsection}
\providecommand{\nobreakdash}{\nobreak\hbox{-}}
\setlength{\emergencystretch}{2em}
\multlinegap0pt
\usepackage{bm}
\usepackage{bbm}
\usepackage{dsfont}
\usepackage{xspace}
\usepackage{cancel}
\usepackage{mathtools}
\usepackage{amsthm}
\makeatletter
\@ifundefined{theorem}{\newtheorem{theorem}{Theorem}}{}
\@ifundefined{lem}{\newtheorem{lem}[theorem]{Lemma}}{}
\makeatother
\title[On three-dimensional Poisson quasi-Nijenhuis manifolds and H]{On three-dimensional Poisson quasi-Nijenhuis manifolds and Haantjes structures}
\author{E. Chu\~no Vizarreta, I. Mencattini, M. Pedroni}
\date{Source: arXiv:2502.16559v3 (CC BY 4.0)}
\begin{document}
\maketitle

\noindent\emph{Source and licence.} On three-dimensional Poisson quasi-Nijenhuis manifolds and Haantjes structures. Version arXiv:2502.16559v3 (CC BY 4.0), distributed under the Creative Commons Attribution 4.0 International licence, as recorded in the arXiv licence metadata for this exact version and shown on the abstract page https://arxiv.org/abs/2502.16559v3. Full rights notice: licenses/arxiv-2502.16559-CC-BY-4.0.txt.\par\smallskip

\noindent\textit{Editorial scope.} Counted unit: the Lem whose source label is \texttt{lem:phik} in the article (source0.tex lines 1098--1104, source label \texttt{lem:phik}), delivered together with the article's own complete original proof of it (source0.tex lines 1105--1173, one \texttt{proof} environment of 69 lines). No mathematics is written, repaired or re-derived: statement and proof are the article's own text line for line. Required context retained uncounted: eq:phis (lines 321--323); theo:charN (lines 414--428); eq: T\_N 3D PqN (lines 487--489). Every definition, display and label of those lines is retained unchanged. Omitted from this package and not redistributed: 1--320, 324--413, 429--486, 490--1097, 1174--1421. Nothing inside a retained excerpt is omitted or altered: every retained body is delivered exactly as the article's lines are written, so no omission receipt of any kind accompanies this package. Citations: the retained excerpts carry no citation command at all, so no bibliography entry of the article is transcribed and no reference is redistributed. Reference adaptation: none was needed and none was made; every reference inside the delivered unit resolves inside it, so no unresolved reference or citation is produced. Printed slips kept literal: no printed word, symbol, hypothesis or number of the counted statement or of its proof was altered. Layout and class: the article's own class is \texttt{article} with packages including amssymb, amsmath, amsthm, calc, stackrel, color, cancel, lineno, hyperref, authblk, graphicx, comment; the wrapper is the lane's pinned \texttt{amsart} editorial wrapper at 10pt on A4, which restates the theorem-like environments the retained text needs and adds the article's own macro definitions that the retained text needs (none), with extra wrapper packages (bm, bbm, dsfont, xspace, cancel, mathtools). The delivered numbering is the unit's own. Assets: no retained span contains a figure environment, a graphics-inclusion command, a PDF-page inclusion, a table or a TikZ picture; no image file is copied into this package, the package declares no asset dependency and no image file is redistributed with it. Licence: version arXiv:2502.16559v3 of this article is distributed under the Creative Commons Attribution 4.0 International licence, as recorded in the arXiv licence metadata for this exact version and shown on the abstract page https://arxiv.org/abs/2502.16559v3. A full rights notice is delivered at licenses/arxiv-2502.16559-CC-BY-4.0.txt. Characters: every retained excerpt is pure ASCII, so the delivered engine, XeLaTeX with its default Latin Modern text font, reports no missing character.
\begin{equation}
\langle\phi_s,X\rangle=\frac{1}{2}\text{Tr}(N^s(i_XT_N)),\qquad\forall X\in\mathfrak X(M), s\geq 0, \label{eq:phis}
\end{equation}

\begin{theorem}\label{theo:charN}
Let $N$ be a $(1,1)$ tensor field on a three-dimensional oriented Poisson manifold $(M,\pi,\mathcal V)$. Suppose that $\pi$ never vanishes.
Then $(M,\pi,N)$ is a PN manifold if and only if $N$ can be decomposed as 
\begin{equation}
N=\lambda I+Z\otimes\xi,\label{eq:chPN}
\end{equation}
where $\lambda\in C^\infty(M)$, $Z\in\mathfrak X(M)$, $\xi=i_\pi\mathcal V$ and the following conditions hold:
\begin{equation}
d\left(\lambda+\left<\xi,Z\right>\right)=\operatorname{div}(Z)\xi\quad\text{and}\quad Z(\lambda)=0.
\label{eq:contor}
\end{equation}
%\begin{equation}
%\xi\wedge(d\left(\left<\xi,Z\right>+\lambda\right))=0\quad\text{and}\quad Z(\lambda)=0.\label{eq:contor}
%\end{equation}
\end{theorem}

\begin{equation}
T_N(X,Y)=Z(\lambda)i_\xi(X\wedge Y).\label{eq: T_N 3D PqN}
\end{equation}

\begin{lem}\label{lem:phik}
	On an oriented three-dimensional PqN manifold, where $N=\lambda I+Z\otimes\xi$, we have that 
	\begin{equation}
		\phi_k=Z(\lambda)\lambda^k\xi\qquad\mbox{for all $k\geq 0$.}
		\label{eq:phik}
	\end{equation}
\end{lem}

\begin{proof}
	The first step is to prove %(see Section \ref{subsec:app2}) 
	that 
	\begin{equation}
		N^k=\lambda^kI+f_k \,Z\otimes\xi,\label{eq:Nalk}
	\end{equation}
	where $f_k=\sum_{l=0}^{k-1}{k\choose l}\lambda^l\langle\xi,Z\rangle^{k-l-1}$. 
The proof is done by induction, noticing that for $k=1$ one has $N=\lambda I+Z\otimes \xi$, which is the content of Theorem \ref{theo:charN}.
Suppose now that \eqref{eq:Nalk} holds for $k$ and let us prove it for $k+1$. We compute
\begin{equation}
\label{eq:s1}
\begin{aligned}
N^{k+1}(X)&=(\lambda I+Z\otimes\xi)\Big(\lambda^kX+\sum_{l=0}^{k-1}{k\choose l}\lambda^l\langle\xi,Z\rangle^{k-l-1}\langle\xi,X\rangle Z\Big)
%\nonumber
\\
&=\lambda^{k+1}X+\underbrace{\sum_{l=0}^{k-1}{k\choose l}\lambda^{l+1}\langle\xi,Z\rangle^{k-l-1}\langle\xi,X\rangle Z}_{(1)}+\underbrace{\sum_{l=0}^{k-1}{k\choose l}\lambda^{l}\langle\xi,Z\rangle^{k-l}\langle\xi,X\rangle Z+\lambda^k\langle\xi,X\rangle Z}_{(2)}.
\end{aligned}
\end{equation}
Now we note that in \eqref{eq:s1} one has
\begin{eqnarray}
&(1)&\stackrel{s=l+1}{=}\sum_{s=1}^{k}{k\choose s-1}\lambda^{s}\langle\xi,Z\rangle^{k-s}\langle\xi,X\rangle Z,\label{eq:s2}\\
&(2)&=\sum_{l=0}^k{k\choose l}\lambda^l\langle\xi,Z\rangle^{k-l}\langle\xi,X\rangle Z=\sum_{l=1}^k{k\choose l}\lambda^l\langle\xi,Z\rangle^{k-l}\langle\xi,X\rangle Z+\langle\xi,Z\rangle^k\langle\xi,Z\rangle Z.\label{eq:s3}
\end{eqnarray}
By renaming $l$ the running variable $s$ in \eqref{eq:s2}, plugging \eqref{eq:s2} and \eqref{eq:s3} in \eqref{eq:s1}, 
and using the identity ${k\choose l-1}+{k\choose l}={k+1\choose l}$, one obtains 
\begin{eqnarray*}
%&&
N^{k+1}(X)%\\
&=&\lambda^{k+1}X+\sum_{l=1}^{k}\Big[{k\choose l-1}+{k\choose l}\Big]\lambda^{l}\langle\xi,Z\rangle^{k-l}\langle\xi,X\rangle Z+\langle\xi,Z\rangle^k\langle\xi,Z\rangle Z\\
&%\stackrel{{k\choose l-1}+{k\choose l}={k+1\choose l}}{=}
=&\lambda^{k+1}X+\sum_{l=1}^{k}{k+1\choose l}\lambda^{l}\langle\xi,Z\rangle^{k-l}\langle\xi,X\rangle Z+\langle\xi,Z\rangle^k\langle\xi,Z\rangle Z\\
&=&\lambda^{k+1}X+\sum_{l=0}^{k}{k+1\choose l}\lambda^{l}\langle\xi,Z\rangle^{k-l}\langle\xi,X\rangle Z,
\end{eqnarray*}
which, being $X$ arbitrary, entails the identity \eqref{eq:Nalk}.  

The second step starts with the remark that, since
	\[
	T_N(X,Y)\stackrel{\eqref{eq: T_N 3D PqN}}{=}Z(\lambda)i_\xi(X\wedge Y)=Z(\lambda)\langle\xi,X\rangle Y-Z(\lambda)\langle\xi,Y\rangle X,
	\]
	one has that
	\begin{equation}
		i_XT_N=Z(\lambda)(\langle\xi,X\rangle I-X\otimes\xi).\label{eq:iXTN}
	\end{equation}
	For every $k\geq 0$ and $X\in\mathfrak X(M)$, we can now compute
	\[
	\begin{split}
		\langle\phi_k,X\rangle&\stackrel{%\eqref{eq:phi}
			\eqref{eq:phis}}{=}\frac12\text{Tr}(N^ki_XT_N)\stackrel{\eqref{eq:iXTN}}
		{=}\frac12\text{Tr}\big(N^k(Z(\lambda)\langle\xi,X\rangle I-Z(\lambda)X\otimes\xi)\big)\\
		&\stackrel{\eqref{eq:Nalk}}{=}\frac12\text{Tr}\Big[\Big(\lambda^kI+\sum_{l=0}^{k-1}{k\choose l}\lambda^l\langle\xi,Z\rangle^{k-l-1}Z\otimes\xi \Big)\Big(Z(\lambda)\langle\xi,X\rangle I-Z(\lambda)X\otimes\xi)\Big)\Big]\\
		&=\frac12\text{Tr}\Big[Z(\lambda)\lambda^k\langle\xi,X\rangle I-Z(\lambda)\lambda^kX\otimes\xi+Z(\lambda)\langle\xi,X\rangle\sum_{l=0}^{k-1}{k\choose l}\lambda^l\langle\xi,Z\rangle^{k-l-1}Z\otimes\xi\\
		&-Z(\lambda)\sum_{l=1}^{k-1}{k\choose l}\lambda^l\langle\xi,Z\rangle^{k-l-1}\langle\xi,X\rangle Z\otimes\xi\Big],
	\end{split}
	\]
	where, in the last sum, $\langle\xi,X\rangle Z\otimes\xi$ comes from the  computation
	\[
	\left((Z\otimes\xi)\circ (X\otimes\xi)\right)(Y)=\langle\xi,Y\rangle(Z\otimes\xi)(X)=\langle\xi,Y\rangle\langle\xi,X\rangle Z%,
	\qquad\forall Y\in\mathfrak X(M),
	\]
	which yields the identity $(Z\otimes\xi)\circ(X\otimes\xi)=\langle\xi,X\rangle Z\otimes\xi$. Now, recalling that $\text{Tr}(Z\otimes\xi)=\langle\xi,Z\rangle$ and $\text{Tr}(I)=3$, the last term of the previous chain of identities becomes
	\[
	\begin{split}
		&\frac12\left(3Z(\lambda)\lambda^k\langle\xi,X\rangle-\lambda^kZ(\lambda)\langle\xi,X\rangle+\cancel{Z(\lambda)\langle\xi,X\rangle\sum_{l=0}^{k-1}{k\choose l}\lambda^l\langle\xi,Z\rangle^{k-l}}\cancel{-Z(\lambda)\langle\xi,X\rangle\sum_{l=0}^{k-1}{k\choose l}\lambda^l\langle\xi,Z\rangle^{k-l}}
		\right)\\
		&=Z(\lambda)\lambda^k\langle\xi,X\rangle,
	\end{split}
	\]
	which entails $\phi_k=Z(\lambda)\lambda^k\xi$.
\end{proof}

\end{document}
